% ##############################################################################
\section{Safe View-As Without Shared Accounts}
\label{sec:impersonation}

Support engineers often need to see a tenant's UI exactly as that tenant sees it. The lazy answer is a shared admin account. It works until it matters. Then the audit trail points to a credential instead of a person, the blast radius is concentrated in one login, and every incident review starts with a shrug.

\system{} uses a view-as workflow instead. The actor stays fixed. The effective context changes.

\subsection{Actor and effective context}
\label{sec:actor_context}

The \code{actor} field in \code{/v2/me} is the authenticated identity. It does not change during a session. The \code{view_as} field is the adopted tenant and role. Every request made during a support session carries both values in logs, so the system can answer two different questions.

\begin{itemize}
  \item Who actually clicked the button?
  \item Which tenant and role were they viewing when they clicked it?
\end{itemize}

Those questions should never share one field.

\subsection{Server-enforced allow-lists}
\label{sec:allowlists}

The control plane computes the \code{impersonatable} list from the actor's role and support scope. The web tier may display this list, but it does not own it. When a support engineer activates view-as, the control plane validates the requested tenant and role against the server-computed allow-list before applying the effective context.

A modified browser can ask for another tenant. The middleware still says no.

\subsection{Signed view-as headers}
\label{sec:signed_headers}

The web tier signs the effective context with HMAC-SHA256. The signed payload contains the tenant, role, and a per-activation nonce. The signing key stays server-side; the browser never receives it. The nonce is registered at activation and revoked on logout or tenant switch. If nonce registration fails, activation fails closed.

\Needspace{10\baselineskip}
\begin{lstlisting}[style=typescript, caption={Sanitized view-as activation sketch.}, label={lst:viewas_sign}]
const nonce = randomUUID();
const viewAs = `tenant=${tenantId};role=${role};nonce=${nonce}`;
const signature = signWithServerKey(viewAs);
\end{lstlisting}

The signature is not the only control. API Gateway still validates the Cognito token, the Lambda still checks the allow-list, backend APIs still enforce authorization, and logs still contain immutable actor attribution. The signing layer narrows replay and header-injection risk; it does not carry the whole security model.
